\chapter{Obstruction atlases and arithmetic activation}\label{ch:atlas}

Chapter~\ref{ch:topology} analysed one explicit cycle. The census of \cref{obs:census-H2} shows that this cycle explains only part of the $H_2$ that arithmetic produces. This chapter replaces the single cycle by a \emph{catalogue}: for any cycle $z$ of the permutohedron the survival law has the same form as for the anchor, and the family of laws that occur is itself an object of study. The theorem of \cref{sec:atlas-law} is general; the catalogue of \cref{sec:atlas-census} is a finite census with an explicit prediction record and an explicit statement of what it does not show.

\section{The survival law of an arbitrary cycle}\label{sec:atlas-law}
Fix $k$ and write $A(m)=\mathfrak F_\lambda(m)\subseteq 2^{[k]}\setminus\{\emptyset\}$ for the family of feasible prefix sets of a content at source $m$ (\cref{lem:F}). For a $j$-cycle $z$ of $P_k$ over $\F$ put $P_z=\bigcup_{F\in\operatorname{supp}z}\Pref(F)$ and $K_z(\lambda)=\lcm_{I\in P_z}D(I)$.
\begin{definition}[minimal filling family]\label{def:filling-family}
For a $j$-cycle $z$ of $P_k$ call $S\subseteq 2^{[k]}\setminus P_z$ \emph{filling} if $z\in\operatorname{im}\partial_{j+1}$ using only those $(j{+}1)$-faces $F$ of $P_k$ with $\Pref(F)\subseteq P_z\cup S$. The \emph{minimal filling families} $\mathfrak M(z)$ are the inclusion-minimal filling sets.
\end{definition}
Filling is monotone: enlarging $S$ enlarges the set of available $(j{+}1)$-faces, hence the image of $\partial_{j+1}$. So the filling sets form an up-set, determined by the antichain $\mathfrak M(z)$.
\begin{theorem}[survival law for a given cycle]\label{thm:AT1}
Let $z$ be a $j$-cycle of $P_k$. For every content $\lambda$ and source $m$,
\begin{gather*}
z\subseteq X_\lambda(m)\ \text{ and }\ [z]\ne0\in H_j(X_\lambda(m);\F)\\
\iff\ P_z\subseteq A(m)\ \text{ and no } S\in\mathfrak M(z)\text{ satisfies } S\subseteq A(m).
\end{gather*}
Consequently (existence): a source $m$ with $[z]\neq0$ exists iff no $S\in\mathfrak M(z)$ satisfies $S\subseteq A(K_z(\lambda))$, and then $m=K_z(\lambda)$ is a witness.
\status{\Analytic\ (the law and the existence criterion for every $j$ and $k$); the lists $\mathfrak M(z)$ themselves are \FiniteExact\ for the cycles of \cref{obs:atlas} and \CompAss\ in general}{GIG-I v005 Stage 2, AT-1, AT-2; v005 Stage 3, AT-2$'$}
\end{theorem}
\begin{proof}
By \cref{thm:face-threshold}, a face $F$ lies in $X_\lambda(m)$ iff $\Pref(F)\subseteq A(m)$. Hence $z\subseteq X$ iff $P_z\subseteq A(m)$. Given that, $[z]=0$ in $H_j(X)$ iff $z=\partial_{j+1}c$ for a chain $c$ supported on $(j{+}1)$-faces of $X$, i.e. with $\Pref\subseteq A(m)$; this is possible iff $A(m)\setminus P_z$ contains a filling set (the faces used have $\Pref(F)\subseteq P_z\cup S$ with $S=A(m)\setminus P_z$, and conversely a filling set contained in $A(m)\setminus P_z$ makes its faces present), iff $A(m)$ contains some $S\in\mathfrak M(z)$ (up-set property). Negating gives the law.
For existence: $P_z\subseteq A(m)$ iff $K_z\mid m$ (\cref{lem:F}); and $m\mid m'$ implies $A(m)\subseteq A(m')$ (\cref{thm:face-threshold}(b)). If $[z]\ne0$ at some $m$, then $K_z\mid m$, so $A(K_z)\subseteq A(m)$, and no $S\in\mathfrak M(z)$ lies in $A(m)\supseteq A(K_z)$. Conversely $P_z\subseteq A(K_z)$, so survival holds at $m=K_z$ when no $S$ lies in $A(K_z)$.
\end{proof}
\begin{remark}[what the theorem says and does not say]
The theorem is a restatement of the filtration: nothing arithmetic enters beyond the monotonicity of $A(m)$ in $m$ under divisibility. Its content is that \emph{every} survival question reduces to a monotone Boolean function of the $\le 2^k$ prefix sets, whose minimal true points $\mathfrak M(z)$ are a finite combinatorial invariant of $(P_k,z)$ alone. Arithmetic then enters only through which sets lie in $A(m)$ (\cref{prop:valuation}), and \cref{prop:CRT} says how freely these can be prescribed. For the anchor, $\mathfrak M(z)=\{B_0,B_1,B_2,B_3\}$ (\cref{thm:exact-survival}); \cref{thm:hole-survival} is the weaker statement that follows from the fact that each $B_i$ contains a gate.
\end{remark}

\subsection{Computing \texorpdfstring{$\mathfrak M(z)$}{M(z)}}
The families were computed in the source programme by \emph{dualisation}: the set of non-filling families is a down-set whose maximal elements are the \emph{maximal non-filling} sets $N_1,\dots,N_t$; a set $S$ fills iff it meets the complement of every $N_i$, so $\mathfrak M(z)$ is the set of minimal transversals of $\{\overline{N_i}\}$. The algorithm alternates Berge dualisation with greedy extension of a non-filling transversal to a maximal non-filling set, and stops when every minimal transversal tests as filling; then each is a minimal true point and every true point contains one. This is exact (no enumeration cut-off), the completeness certificate being independently checked. \Cref{app:methods} records which of these computations the manuscript verifier re-derived (the anchor, by direct testing against all $2^{20}$ candidates restricted to a sample, \cref{thm:exact-survival}) and which are inherited.

\section{The atlas}\label{sec:atlas-census}
\begin{observation}[the obstruction atlas]\label{obs:atlas}
The $158$ hole regimes of the six-label census (\cref{obs:census-H2}) carry $160$ homology-basis cycles. Reduced to canonical form (minimal weight representatives, modulo the action of $S_6$ on ranks) they fall into $70$ canonical cycles and these into $23$ \emph{law classes}: classes of pairs $(P_z,\mathfrak M(z))$ under $S_6$. Nine further classes arise from held-out failures (below), giving $32$ classes. Of the $32$: $24$ have $4$ minimal families, $7$ have $6$ and one has $8$; the number of singleton families (single necessary prefix sets, ``universal gates'') is $0$ for $8$ classes, $1$ for $10$ and $2$ for $14$. The anchor of \cref{ch:topology} is class~$0$; it alone covers $59$ regimes (seven more than \cref{thm:hole-survival} explains) and leaves $99$ of the $106$ ``unexplained'' regimes uncovered, while the $32$ classes together cover all $158$.
\status{\FiniteExact\ as a statement about the listed cycles (inherited, \texttt{NOT\_RERUN}, certificates named in \cref{app:sources}); coverage \Bounded}{GIG-I v005 Stage 2, AT-3; E22--E26}
\end{observation}
\begin{observation}[prediction record]\label{obs:atlas-pred}
By \cref{thm:AT1} a class predicts a hole regime iff its law holds there, so a class can never produce a false positive. The atlas was tested on contents with at least one label above $16$ (disjoint from the census), computing exact $\beta_2$ at every regime. \Cref{tab:atlas-pred} lists the three rounds.
\begin{table}[ht]\centering\small
\begin{tabular}{@{}lrrrrr@{}}
\toprule
Round & contents & regimes & hole regimes & missed & false pos.\\ \midrule
Atlas v1 (23 classes), set 1 & 600 & 33{,}584 & 211 & 18 & 0\\
Atlas v1, set 2 & 600 & -- & 171 & 13 & 0\\
Atlas v2 (+18 from set 1), set 2 & 600 & -- & 171 & 0 & 0\\
Atlas v2, set 3 & 1000 & 54{,}091 & 292 & 3 & 0\\ \bottomrule
\end{tabular}
\caption{Held-out prediction of $H_2$ regimes by the atlas, labels $\le40$. ``Hole regimes'' is the number with $\beta_2>0$ (set~1: $193$ predicted $+\,18$ missed; set~3: $289+3$); ``missed'' counts false negatives. A dash marks a figure not recorded in the source summary. Atlas v2 was updated only with the false negatives of set~1 and then scored unchanged on sets~2--3.}\label{tab:atlas-pred}
\end{table}
New classes still appear at about $1\%$ of hole regimes for larger labels, so the atlas is \emph{not closed}.
\status{\Bounded\ (inherited, \texttt{NOT\_RERUN}; the zero false-positive record is a consequence of \cref{thm:AT1}, not an empirical accident)}{GIG-I v005 Stage 2, AT-4; E24--E26}
\end{observation}

\section{Limits of the atlas}
We state explicitly what the atlas does not establish.
\begin{enumerate}[label=(\roman*)]
\item \emph{Finiteness.} It is not shown that the number of law classes is finite for six labels. The observed new-class rates ($18/211$, $13/171$, $3/292$ on successive held-out rounds) are decreasing but no bound is proved.
\item \emph{Representative dependence.} Classes are defined from low-weight representatives of homology classes. Different representatives of one homology class can have different laws; whether the $32$ classes merge when defined from the homology class itself (for example through the multi-parameter persistence of $[z]$ over the divisibility lattice, cf.\ \cref{ch:synthesis}) is open.
\item \emph{Arithmetic realisability.} \Cref{prop:gate-nec} shows that the large families of the anchor class can be arithmetically realised or avoided. The analogous question for the other $31$ classes (which families are realisable, and which large families are arithmetically equivalent to smaller ones) is untested.
\item \emph{Independent description.} No construction of the cycles or families independent of the finite census is known; in particular the $50$-, $62$- and $74/76$-face cycle families (two singleton gates and two large families) are described only as data.
\end{enumerate}
These are listed as open problems in \cref{ch:open} (\cref{prob:atlas-closure,prob:representatives,prob:realisability}).
\begin{proposition}[the atlas is a statement about contents, not about $H_2$ in general]\label{prop:atlas-scope}
Every atlas statement is relative to (a) six-label contents, (b) the GF(2) field, (c) the SELECT-II family $\mathfrak F_\lambda(m)=\{I:D(I)\mid m\}$. Replacing (c) by another prefix-defined system leaves \cref{thm:AT1} intact (its proof uses only the prefix-defined property) but changes which families $A(m)$ are attainable and hence which classes occur.
\status{\Analytic}{GIG-I v005 Stage 2, AT-1 (proof)}
\end{proposition}
